\documentclass[11pt,a4paper]{article}
\usepackage[T1]{fontenc}
\usepackage[utf8]{inputenc}
\usepackage[ngerman]{babel}
\usepackage[margin=2.5cm]{geometry}
\usepackage[hidelinks]{hyperref}
\usepackage{array}
\setlength{\parindent}{0pt}
\setlength{\emergencystretch}{2em}
\renewcommand{\arraystretch}{1.25}
\begin{document}
\shorthandoff{"}
\begin{center}
  {\Large\textbf{Webtechnologien -- Übungsblatt 2}}
\end{center}
\section*{Aufgabe 6: Image Gallery}
\textbf{Aufgabe:} Die Bildergalerie aus der Vorlesung verstehen und im
Vollbildmodus um die Buttons \glqq previous\grqq{} und \glqq next\grqq{}
erweitern. Nur neuer Code in \texttt{gallery.js}, bestehender Code bleibt
unverändert.

\medskip
\textbf{HTML-Struktur:}
\begin{itemize}
  \item \texttt{div.gallery} mit den Vorschaubildern. Jedes \texttt{img}
    hat \texttt{src} (kleine Datei), \texttt{data-large} (großes Bild) und
    \texttt{alt} (Beschreibung).
  \item \texttt{div\#displayDIV} mit \texttt{img.bigIMG} und
    \texttt{p.desc}. Laut CSS zuerst unsichtbar (\texttt{display: none}),
    sonst ein silberner Vollbildbereich mit $100\,\mathrm{vw} \times
    100\,\mathrm{vh}$ und \texttt{position: absolute}.
\end{itemize}

\textbf{Originalcode:}
\begin{itemize}
  \item \texttt{showImg(e)}: setzt \texttt{bigIMG.src} auf
    \texttt{e.target.dataset.large}, die Beschreibung auf
    \texttt{e.target.alt} und zeigt \texttt{displayDIV} an.
  \item \texttt{hideImg(e)}: setzt \texttt{e.currentTarget.style.display}
    auf \texttt{none}. Ein Klick irgendwo auf das Vollbild schließt es.
  \item Eine Schleife hängt \texttt{showImg} an alle
    \texttt{.gallery img}, \texttt{hideImg} hängt an \texttt{displayDIV}.
\end{itemize}

\textbf{Ergebnis} (Erweiterung ab Zeile 17):
\begin{itemize}
  \item \texttt{currentIndex} speichert das angezeigte Bild. Ein zweiter
    \texttt{click}-Listener pro Vorschaubild setzt \texttt{currentIndex = i}.
    Durch \texttt{let} in der \texttt{for}-Schleife hat jede Runde ihr
    eigenes \texttt{i}.
  \item \texttt{showImgAt(index)}: berechnet den neuen Index mit
    \begin{center}
    \texttt{currentIndex = (index + images.length) \% images.length}
    \end{center}
    und setzt Bild und Beschreibung wie das Original.
    Bei drei Bildern: Index 0, previous: $(0 - 1 + 3) \bmod 3 = 2$
    (letztes Bild); Index 2, next: $(2 + 1 + 3) \bmod 3 = 0$ (erstes Bild).
  \item \texttt{addButton(text, step, side)}: erzeugt einen
    \texttt{button}, absolut positioniert (\texttt{top: 50\%}, 20\,px vom
    linken bzw. rechten Rand), und hängt ihn in \texttt{\#displayDIV}.
    Der Klick ruft \texttt{showImgAt(currentIndex + step)} auf.
  \item \texttt{e.stopPropagation()} im Button-Listener verhindert, dass der
    Klick zu \texttt{displayDIV} hochwandert und \texttt{hideImg} das Bild
    sofort schließt.
  \item Aufrufe: \texttt{addButton("previous", -1, "left")} und
    \texttt{addButton("next", 1, "right")}.
\end{itemize}
Die Erweiterung verwendet die Liste \texttt{images} aus dem Originalcode,
also alle Bilder in \texttt{.gallery}. Sie funktioniert daher mit
beliebig vielen Bildern im gleichen HTML-Format.

\medskip
\textbf{Test:} In Chrome, \texttt{gallery.html} direkt geöffnet.
Erstes Bild, next: \texttt{street.JPG} mit \glqq Hong Kong central district
life\grqq. Danach zweimal previous: \texttt{taichung.JPG} (Sprung zum
letzten Bild), dann next: \texttt{hk.JPG} (Sprung zum ersten Bild).
Drittes Bild öffnen und next: \texttt{hk.JPG}. Ein Klick neben die
Buttons schließt das Vollbild. Keine Fehler in der Konsole.

\medskip
{\small \textbf{Hilfsmittel:} Claude Code (Claude Opus 5.5) für Erstellung, Prüfung und Dokumentation.}
\end{document}
